\section*{NeurIPS Paper Checklist}

\begin{enumerate}

\item {\bf Claims}
    \item[] Question: Do the main claims made in the abstract and introduction accurately reflect the paper's contributions and scope?
    \item[] Answer: \answerYes{}
    \item[] Justification: The abstract states three claims: (i) rotation-invariant scores collapse to dimension-ranking in heterogeneous pools; (ii) SSR$^{\ddagger}$ halves regret over RankMe; (iii) a two-stage pipeline preserves quality at one-third cost. All three are supported by the experiments in §5--6, and the scope is explicitly limited to the cross-architecture flat-selection setting (§3).

\item {\bf Limitations}
    \item[] Question: Does the paper discuss the limitations of the work performed by the authors?
    \item[] Answer: \answerYes{}
    \item[] Justification: §7 (Discussion) contains an explicit ``Limitations'' paragraph that states: no theoretical bound on LP accuracy in terms of SSR$^{\ddagger}$; MV pool two-stage comparisons have wide CIs at $n{=}5$ tasks; within-family ranking is out of scope; LP ground truth uses a single 80/20 split.

\item {\bf Theory assumptions and proofs}
    \item[] Question: For each theoretical result, does the paper provide the full set of assumptions and a complete (and correct) proof?
    \item[] Answer: \answerYes{}
    \item[] Justification: Lemma~1, Corollary~1, and Proposition~1 in §4 each state all assumptions explicitly (near-isotropy, $n \gg d$, right-rotation invariance). Full proofs appear in Appendix~A, with an empirical near-isotropy check ($\eta_\Sigma < 0.1$ for 93.3\% of pairs) also in Appendix~A.

\item {\bf Experimental result reproducibility}
    \item[] Question: Does the paper fully disclose all the information needed to reproduce the main experimental results of the paper to the extent that it affects the main claims and/or conclusions of the paper (regardless of whether the code and data are provided or not)?
    \item[] Answer: \answerYes{}
    \item[] Justification: All hyperparameters (SA$_k$ dimension $k{=}20$, reference task STL-10, LP regularisation $C{=}1$, 80/20 split, $n{=}1000$ samples per file) are reported in §5 and in the pool appendices (Appendix~B--E). The SSR$^{\ddagger}$ formula is given in closed form in §5. Bootstrap CI details (10\,000 resamples, 95\% level, task-level resampling) are in §5.

\item {\bf Open access to data and code}
    \item[] Question: Does the paper provide open access to the data and code, with sufficient instructions to faithfully reproduce the main experimental results, as described in supplemental material?
    \item[] Answer: \answerYes{}
    \item[] Justification: An anonymised repository (\url{https://anonymous.4open.science/r/crossarchitecturetransferability-FEC8}) ships all scoring / analysis code plus 16 pre-computed result CSVs. A \texttt{reproduce\_all.py} script with \texttt{--verify} flag cross-checks every headline number against the shipped CSVs. Feature extraction scripts and a \texttt{FEATURE\_EXTRACTION.md} guide are included for full re-runs from scratch.

\item {\bf Experimental setting/details}
    \item[] Question: Does the paper specify all the training and test details (e.g., data splits, hyperparameters, how they were chosen, type of optimizer) necessary to understand the results?
    \item[] Answer: \answerYes{}
    \item[] Justification: No model training is involved; the method is a scoring function applied to pre-extracted features. All relevant settings (split ratio, regularisation $C$, sample size $n$, subspace size $k$, reference encoder) are specified in §5 and Appendices~B--E. The choice of $k{=}20$ and STL-10 reference are motivated by ablations in Appendix~D.

\item {\bf Experiment statistical significance}
    \item[] Question: Does the paper report error bars suitably and correctly defined or other appropriate information about the statistical significance of the experiments?
    \item[] Answer: \answerYes{}
    \item[] Justification: All headline regret comparisons report 95\% CIs via 10\,000-resample task-level bootstrap (Tables~2--4). Cross-pool significance: paired bootstrap on Table~7 (23 task units) gives $p < 10^{-4}$ for SSR$^{\ddagger}$ vs.\ RankMe; an additional Wilcoxon signed-rank test on 22 paired units (excluding STL-10 reference) gives $p = 0.0036$. The CI construction and test choice are described in §5. Where $n{=}5$ tasks yield wide CIs (MV pool two-stage), this is explicitly noted as a limitation.

\item {\bf Experiments compute resources}
    \item[] Question: For each experiment, does the paper provide sufficient information on the computer resources (type of compute workers, memory, time of execution) needed to reproduce the experiments?
    \item[] Answer: \answerYes{}
    \item[] Justification: Feature extraction (one-time) requires a GPU (tested on a single V100/A100; $\approx$6 h for all four pools combined). All scoring and analysis runs on CPU in under 10 minutes using the pre-extracted features. This is stated in the repository \texttt{FEATURE\_EXTRACTION.md}. No distributed or proprietary compute is required.

\item {\bf Code of ethics}
    \item[] Question: Does the research conducted in the paper conform, in every respect, with the NeurIPS Code of Ethics?
    \item[] Answer: \answerYes{}
    \item[] Justification: The work is a methodological contribution to label-free model selection using publicly available pretrained models and benchmark datasets. No private data, no human subjects, no sensitive applications, and no data was collected or generated beyond running standard feature extractors on public benchmarks.

\item {\bf Broader impacts}
    \item[] Question: Does the paper discuss both potential positive societal impacts and negative societal impacts of the work performed?
    \item[] Answer: \answerNA{}
    \item[] Justification: This is foundational methodology research on label-free model selection. The method helps practitioners choose among existing pretrained models more accurately, which may reduce wasted compute. There is no direct path to any harmful application; no new training data, proprietary models, or personal data is involved.

\item {\bf Safeguards}
    \item[] Question: Does the paper describe safeguards that have been put in place for responsible release of data or models that have a high risk for misuse (e.g., pre-trained language models, image generators, or scraped datasets)?
    \item[] Answer: \answerNA{}
    \item[] Justification: The paper releases only a scoring algorithm (Python code) and pre-computed numerical result tables. No new pretrained models, datasets, or high-risk generative systems are released.

\item {\bf Licenses for existing assets}
    \item[] Question: Are the creators or original owners of assets (e.g., code, data, models), used in the paper, properly credited and are the license and terms of use explicitly mentioned and properly respected?
    \item[] Answer: \answerYes{}
    \item[] Justification: All pretrained models are drawn from the timm library \citep{wightman2019timm} (Apache 2.0) and torchvision (BSD). Benchmark datasets (CIFAR-10/100, STL-10, DTD, EuroSAT, SVHN, Flowers-102, Pets, Food-101, FGVC-Aircraft, Fashion-MNIST) are publicly available standard benchmarks and are cited. The VTAB benchmark is cited via \citet{zhai2019large}.

\item {\bf New assets}
    \item[] Question: Are new assets introduced in the paper well documented and is the documentation provided alongside the assets?
    \item[] Answer: \answerYes{}
    \item[] Justification: The repository contains (i) the SSR$^{\ddagger}$ scoring code with inline documentation; (ii) pre-computed result CSVs for all 16 reported tables; (iii) a \texttt{README.md} describing every file; (iv) a \texttt{FEATURE\_EXTRACTION.md} with step-by-step GPU setup. All are released under MIT licence in the anonymised repository.

\item {\bf Crowdsourcing and research with human subjects}
    \item[] Question: For crowdsourcing experiments and research with human subjects, does the paper include the full text of instructions given to participants and screenshots, if applicable, as well as details about compensation (if any)?
    \item[] Answer: \answerNA{}
    \item[] Justification: No crowdsourcing or human subjects are involved. All experiments operate on pretrained model features and standard image classification benchmarks.

\item {\bf Institutional review board (IRB) approvals or equivalent for research with human subjects}
    \item[] Question: Does the paper describe potential risks incurred by study participants, whether such risks were disclosed to the subjects, and whether Institutional Review Board (IRB) approvals (or an equivalent approval/review based on the requirements of your country or institution) were obtained?
    \item[] Answer: \answerNA{}
    \item[] Justification: No human subjects are involved in this research.

\item {\bf Declaration of LLM usage}
    \item[] Question: Does the paper describe the usage of LLMs if it is an important, original, or non-standard component of the core methods in this research?
    \item[] Answer: \answerNA{}
    \item[] Justification: LLMs were used only for writing assistance (grammar and phrasing) and do not form any component of the proposed method, experimental design, or evaluation pipeline.

\end{enumerate}
